برای حل این دستگاه معادلات دیفرانسیل جفت‌شده، از خاصیت مشتق‌گیری در حوزه لاپلاس استفاده می‌کنیم. می‌دانیم که $\mathcal{L}\left\{\frac{d}{dt}g(t)\right\} = sG(s) - g(0^-)$. با توجه به این که در صورت سوال ذکر شده سیگنال‌ها راست‌طرف هستند، شرایط اولیه را صفر فرض کرده و از طرفین معادلات تبدیل لاپلاس می‌گیریم:

\begin{align}
sX(s) &= -2Y(s) + 1 \\
sY(s) &= 2X(s) \implies Y(s) = \frac{2}{s}X(s)
\end{align}

حالا رابطه (۲) را در رابطه (۱) جایگذاری می‌کنیم تا $X(s)$ برحسب $s$ به دست آید:
\begin{equation*}
sX(s) = -2\left(\frac{2}{s}X(s)\right) + 1 \implies sX(s) + \frac{4}{s}X(s) = 1
\end{equation*}
\begin{equation*}
X(s)\left( \frac{s^2 + 4}{s} \right) = 1 \implies X(s) = \frac{s}{s^2 + 4}
\end{equation*}

اکنون با استفاده از مقدار $X(s)$ به دست آمده، $Y(s)$ را از روی رابطه (۲) محاسبه می‌کنیم:
\begin{equation*}
Y(s) = \frac{2}{s} \left( \frac{s}{s^2 + 4} \right) = \frac{2}{s^2 + 4}
\end{equation*}

هر دو تابع تبدیل $X(s)$ و $Y(s)$ دارای مخرج مشترک $s^2 + 4 = 0$ هستند. بنابراین قطب‌های سیستم بر روی محور موهومی در نقاط زیر قرار دارند:
\begin{equation*}
s = \pm j2 \implies \text{Re}(s) = 0
\end{equation*}

چون در صورت سوال به صراحت قید شده است که هر دو سیگنال $x(t)$ و $y(t)$ {راست‌طرف} هستند، ناحیه همگرایی برای هر دو تبدیل باید به صورت راست‌بر نسبت به بزرگترین بخش حقیقی قطب‌ها (که در اینجا صفر است) تعریف شود:
\begin{equation*}
\text{ROC}: \text{Re}(s) > 0
\end{equation*}

بنابراین پاسخ نهایی به صورت زیر است:
\begin{align*}
X(s) &= \frac{s}{s^2 + 4}, \quad \text{ROC: } \ \text{Re}(s) > 0 \\
Y(s) &= \frac{2}{s^2 + 4}, \quad \text{ROC: } \ \text{Re}(s) > 0
\end{align*}
